Krátké, ověřené příklady nasazení a doporučené postupy, které lze adaptovat na předmět nebo pilot na úrovni katedry. Níže uvedené příklady vycházejí z praktických doporučení a nástrojů zmíněných ve zdrojovém materiálu; u každého příkladu je uvedena klíčová poznámka k přenositelnosti.

Příklady z praxe
\begin{itemize}
  \item Minecraft Education — výukové světy zaměřené na seznámení se s principy AI (Hour of Code) jsou volně dostupné a škola může použít připravené lekce k demonstraci principů strojového učení a praktických aktivit ve vyučování. Část obsahu je dostupná zdarma; rozšiřující světy jsou součástí školní licence Microsoft 365 A3/A5 \cite{kb_ai_101_tipu_pdf_04e4a115_page_15_2}.
  \item Microsoft výukové nástroje pro matematiku a OneNote — nástroje jako „Pokrok v matematice“ a funkce OneNote pro převod rukopisu na digitální formu usnadňují generování cvičných příkladů, sběr postupů a diagnostiku chyb u studentů; lze je využít pro adaptivní procvičování a sledování pokroku \cite{kb_ai_101_tipu_pdf_04e4a115_page_8_1}.
  \item Profesní rozvoj učitelů — doporučení zaměřit se na vybranou oblast/technologii a učit se v hloubce namísto snahy „stíhat vše“; pro praktické školení lze využít oficiální materiály (např. Microsoft Learn) jako systematický zdroj kurzů a modulků pro pedagogy i IT podporu \cite{kb_ai_101_tipu_pdf_04e4a115_page_54_1}.
\end{itemize}

Klíčové poučení (co si odnést)
\begin{itemize}
  \item Využijte hotové didaktické artefakty (např. připravené světy v Minecraft Education) pro rychlý start s praktickými ukázkami ve výuce \cite{kb_ai_101_tipu_pdf_04e4a115_page_15_2}.
  \item Integrace nástrojů, které zaznamenávají postupy (commit/artefakty, nahrávání postupu), zvyšuje ověřitelnost učení a usnadňuje hodnocení procesu — při výběru nástrojů preferujte školní/enterprise řešení, pokud to právně/technicky dává smysl \cite{kb_ai_101_tipu_pdf_04e4a115_page_15_2,kb_ai_101_tipu_pdf_04e4a115_page_8_1}.
  \item Systematické školení pedagogů a podpůrného personálu (vybraná tematická oblast + oficiální learning moduly) zrychlí adaptaci a sníží riziko nekonzistentního nasazení \cite{kb_ai_101_tipu_pdf_04e4a115_page_54_1}.
\end{itemize}

Rychlý plán zavedení vybraného vzoru (1–2měsíční pilot)
\begin{enumerate}
  \item Vyberte jeden kurz nebo modul a konkrétní didaktický cíl (např. demonstrace principů AI v 90minutové hodině pomocí Minecraft Education). Podle doporučení začněte s úzkou oblastí \cite{kb_ai_101_tipu_pdf_04e4a115_page_54_1}.
  \item Ověřte licenční a datovou situaci (je-li potřeba školní licence Microsoft 365 A3/A5 pro rozšířené světy nebo enterprise řešení pro zpracování studentských dat) a zajistěte DPA / souhlas tam, kde je to nutné \cite{kb_ai_101_tipu_pdf_04e4a115_page_15_2}.
  \item Připravte konkrétní hodinu/aktivitu: stáhněte volnou lekci Hour of Code nebo připravený svět, upravte ji pro lokální kontext, připravte krátký doprovodný materiál pro studenty (cíle, očekávané výstupy) \cite{kb_ai_101_tipu_pdf_04e4a115_page_15_2}.
  \item Nasazení a sběr dat o průběhu: zaznamenávejte artefakty studentské práce a krátkou reflexi (pomůže při vyhodnocení a integritě). Pokud používáte nástroje pro procvičování (math tools, OneNote), zapojte požadavek na nahrání postupu řešení \cite{kb_ai_101_tipu_pdf_04e4a115_page_8_1}.
  \item Vyhodnocení pilotu a škálování: zhodnoťte dosažení cílů, zkušenosti učitele i studentů a rozhodněte o rozšíření (úpravy obsahu, školení dalších učitelů pomocí dostupných modulů, např. Microsoft Learn) \cite{kb_ai_101_tipu_pdf_04e4a115_page_54_1}.
\end{enumerate}

General background knowledge (stručně)
\begin{itemize}
  \item General background knowledge: Ve veřejně známých případech existují širší institucionální nasazení AI (např. Georgia Tech, Harvard CS50, Arizona State, UCL, TU Delft) — tyto kazuistiky slouží jako inspirace pro design patterny a škálování, avšak konkrétní podrobnosti a výsledky je nutné konzultovat v originálních zdrojích (uváděno jako všeobecná reference, není detailně citováno v tomto zdroji).
\end{itemize}